\begin{theorem}[Theorem 1]
Let $n\ge2$, $p>1$ and $r>0$, suppose the Sobolev--Poincar\'e inequality on balls
holds in dimension $n$ (\noderef{CitedSobolevPoincareBall}), and assume the cited
solvability result (C2) (\noderef{CitedDirichletW1p}) and the cited regularity
result (C1) (\noderef{CitedRegularityEpsPos}) in the universal forms described
below.  Then there is a constant $C=C(p,r,n)>0$ with the
following property.  Let $\varepsilon\ge0$, $\gamma_0>0$ and
$\gamma_1\le r\gamma_0$, let $\Omega\subseteq\R^n$ be open, let
$A\colon\R^n\to\R^n$ satisfy $A(0)=0$ and (Q1) of the thesis--(Q2) of the thesis on
$\Omega$ with parameters $p,\varepsilon,\gamma_0,\gamma_1$
(\noderef{EllipticStructureQ12}),
let $u$ be a weak solution of $\operatorname{div}A(\nabla u)=0$ in $\Omega$
(\noderef{IsWeakSolutionOn}) with weak gradient $g=\nabla u$, and let
$B_R=B_R(x_0)$ with $R>0$ and $\overline B_R(x_0)\subseteq\Omega$.  Then
\[
\sup_{B_{R/2}}(\varepsilon+|\nabla u|^2)^{\frac p2}
\;\le\;\frac{C}{R^n}\int_{B_R}(\varepsilon+|\nabla u|^2)^{\frac p2}\,dx .
\]
The supremum is the essential one (p.~6 of the paper), so the conclusion is
written as: for almost every $x\in B_{R/2}$,
\[
(\varepsilon+|g(x)|^2)^{\frac p2}
\le\frac{C}{R^n}\int_{B_R}(\varepsilon+|g|^2)^{\frac p2}\,dy .
\]
The containment hypothesis is $\overline B_R(x_0)\subseteq\Omega$, the closed
ball, as the paper's Note prescribes for Theorems~1 and~2 (``the balls are
concentric and $\overline B_R\subset\Omega$''); this is what makes
$\overline B_R$ a compact subset of $\Omega$ and hence the right-hand side
finite, since \noderef{IsWeakSolutionOn} puts $|g|$ in $L^p$ of the compact
subsets of $\Omega$ only.

The normalisation $A(0)=0$ is \S3's own standing hypothesis, fixed at the start
of the section: ``Since $\operatorname{div}A(0)=0$ we may assume without loss of
generality that $A(0)=0$ and that $A$ satisfies conditions (Q1) of the thesis and
(Q2) of the thesis'' (thesis p.~14).  It is what the cited solvability result (C2)
requires, and it therefore reaches this theorem through the reduction of the
case $\varepsilon=0$ to $\varepsilon>0$, which goes through
\noderef{DegenerateLimitTransfer}; the case $\varepsilon>0$ does not use it.
Replacing $A$ by $A-A(0)$ changes neither (Q1) of the thesis--(Q2) of the thesis, which
involve $A$ only through its derivative, nor the equation, which pairs $A$ with
the gradient of a compactly supported test function; so the hypothesis costs the
statement no generality, which is what the thesis means by ``without loss of
generality''.

The constant $C$ is bound before $\varepsilon,\gamma_0,\gamma_1$, the operator,
the solution and the ball, and the structure constants enter only through the
bound $\gamma_1\le r\gamma_0$: this is the formal content of ``$C$ depends only
on $p$, $\gamma_1/\gamma_0$ and $n$''.

The cited solvability result (C2) (\noderef{CitedDirichletW1p}) is assumed here
for \emph{every} field $A'\colon\R^n\to\R^n$ and all parameters
$\varepsilon',\gamma_0',\gamma_1'$ \emph{with $\gamma_0'>0$ and $\gamma_1'>0$},
which is the form (C2) has in the paper's Note.  The hypothesis is bound before the constant, but it introduces no
dependence on any quantity outside the list the paper names: it mentions only
$n$ and $p$, both already bound, and it carries no numerical content of its own.
It is needed because the reduction of the case $\varepsilon=0$ to
$\varepsilon>0$ goes through \noderef{DegenerateLimitTransfer}, whose own (C2)
hypothesis is of that universal shape --- it has to be, since the regularized
fields to which (C2) is applied are produced inside that node's conclusion.

The cited regularity result (C1) (\noderef{CitedRegularityEpsPos}) is assumed in
the same way: for \emph{every} field $A'$, \emph{every} open set $\Omega'$ and
all parameters $\varepsilon',\gamma_0',\gamma_1'$ with
$\gamma_0'>0$ and $\gamma_1'>0$.  That is the form (C1) has in the
paper's Note, which states it for an arbitrary operator and an arbitrary
solution: ``If $\varepsilon>0$ and $A\colon\R^n\to\R^n$ satisfies
(Q1) of the thesis and (Q2) of the thesis, then every weak solution
$u\in W^{1,p}_{\mathrm{loc}}(\Omega)$ of (3.1) of the thesis satisfies
$u\in W^{2,2}_{\mathrm{loc}}(\Omega)$ and $\nabla u\in C(\Omega)$'' (Note for the
formalization, (C1)).  So the universal clause is the faithful transcription and
a clause attached only to the field $A$ and the domain $\Omega$ of this theorem
would be the weakening.  The clause is quantified over
$\varepsilon',\gamma_0',\gamma_1'$ with $\gamma_0'>0$, $\gamma_1'>0$ and over
open sets $\Omega'$ only, and those side conditions are not optional: the
structure conditions (Q1) of the thesis--(Q2) of the thesis carry positive constants in
the thesis (``$\gamma_0,\gamma_1>0$''), and (C1) is a statement about weak
solutions on an open set, so a clause quantified over $\gamma_0'\le0$ or over
arbitrary $\Omega'$ would assert more than the cited result does --- and since
this is a hypothesis of the present theorem, asserting more than the literature
gives risks making the theorem vacuously true.  The remaining antecedents need
no guard, because \noderef{CitedRegularityEpsPos} carries $\varepsilon>0$,
(Q1) of the thesis--(Q2) of the thesis and the weak-solution property inside its own
statement.  The same two positivity conditions are attached to the (C2) clause;
(C2) needs no openness condition, since \noderef{CitedDirichletW1p} is a
statement about balls and fixes its own domain.  Quantifying the set is not decorative: the proof consumes (C1) on
the ball $B_R$, where \noderef{CaccioppoliPowerWeight} needs the second
derivatives, and \noderef{CaccioppoliPowerWeight} takes the family
$D(\nabla u)$ as a parameter, so the proof has to produce one --- it is supplied
by the existential inside (C1)'s conclusion, and there is nothing in this
theorem's own statement for a bare $W^{2,2}_{\mathrm{loc}}$ hypothesis to attach
to.  Like (C2) and (C6), (C1) carries no numerical content: it asserts
membership and continuity and names no constant, so although it is bound before
the existential for $C$ it introduces no dependence of $C$ on anything outside
the list $p$, $\gamma_1/\gamma_0$, $n$ that the paper names.

The hypothesis \noderef{CitedSobolevPoincareBall} is the cited result (C6) of the
paper's Note, carried here as an explicit argument because the proof uses it: the
Moser iteration runs on the scale-invariant Sobolev inequality of
\noderef{SobolevExponentTwoBall}, whose own statement is conditional on (C6).
It too is bound before the existential for $C$ and introduces no dependence of
$C$ on anything outside the paper's list, since (C6) is a statement about $n$
alone and $n$ is already bound first.
\end{theorem}
\begin{proof}
Throughout write $B_\rho(y)$ for the open ball of radius $\rho$ about $y$,
$\overline B_\rho(y)$ for the closed one, $B_\rho=B_\rho(x_0)$ for the balls
concentric with the one in the statement, and $|E|$ for Lebesgue measure.  $D_i$
is the derivative in the direction of the $i$-th standard basis vector,
$(\cdot,\cdot)$ and $|\cdot|$ are the Euclidean inner product and norm, and all
non-integer powers are real powers of nonnegative numbers.

\emph{Step 0: the constants, chosen before all the data.}  Apply
\noderef{DegenerateLimitTransfer} with the present $n,p,r$; its first three
hypotheses $2\le n$, $1<p$, $0<r$ are hypotheses of the present theorem, and it
supplies a number $r'>0$ depending only on $n,p,r$.  Put
\[
\bar r=\max\{r,r'\},
\]
so $\bar r>0$ and $\bar r$ depends only on $n,p,r$.  Let $c>0$ be the constant
that \noderef{CaccioppoliPowerWeight} supplies for the triple $(n,p,\bar r)$ ---
its hypotheses there are $2\le n$, $1<p$ and $\bar r>0$ --- so $c=c(n,p,\bar r)$,
and let $C_0>0$ be the constant that \noderef{MoserLocalBoundedness} supplies for
the pair $(n,Q)$ with $Q=c$; its hypotheses there are $2\le n$, $0<Q$ and the
Sobolev--Poincar\'e inequality on balls in dimension $n$
(\noderef{CitedSobolevPoincareBall}), which is a hypothesis of the present
theorem.  So $C_0=C_0(n,c)$.  Finally put
\[
\kappa=\max\{1,2^{\frac p2-1}\},\qquad
C_1=C_0^2,\qquad
C=C_1\,\kappa^2\,2^{\frac p2} .
\]
All of $r',\bar r,c,C_0,\kappa,C_1,C$ are determined by $n$, $p$ and $r$ alone
and are fixed before $\varepsilon,\gamma_0,\gamma_1,\Omega,A,u,g,x_0,R$, which is
the quantifier order the statement requires; and $C\ge C_1>0$ because
$\kappa\ge1$ and $2^{p/2}\ge1$.  This is the formal content of ``$C$ depends only
on $p$, $\gamma_1/\gamma_0$ and $n$'': the structure constants enter only through
the bound $\gamma_1\le r\gamma_0$.

\emph{Step 1: $\gamma_0\le\gamma_1$, so $\gamma_1>0$.}  This is needed because
the cited results (C1) and (C2) are assumed here only for positive structure
constants.  The set $\Omega$ contains $\overline B_R$ with $R>0$, so it has
positive measure (Preamble: the volume of a ball, so
$|\Omega|\ge|B_R|=\omega_nR^n>0$); by \noderef{EllipticStructureQ12} the
conditions (Q1) of the thesis and (Q2) of the thesis hold for almost every $x\in\Omega$,
i.e.\ each fails only on a null set.  The union of those two null sets is null,
and a null set cannot contain $\Omega$, whose measure is positive; so
$\Omega$ minus that union is nonempty and there is a single $x\in\Omega$ at
which both conditions hold.  \emph{This is the only promotion of an
almost-everywhere clause in this proof that is used at a point}, and it is made
inline here rather than by citation, which is what pm-0034's situation permits
for a field that does not depend on $x$; B1 below makes the same argument once
more, for the field at $\varepsilon=0$.  Fix such an $x$, take $z=0$,
$h=e_1$ --- so that $\varepsilon+|h|^2=\varepsilon+1>0$ and in particular
$h\in D(\varepsilon)$, where the two conditions are asserted --- and $\xi=e_1$.
Then (Q1) of the thesis and (Q2) of the thesis give
\[
\gamma_0(\varepsilon+1)^{\frac p2-1}
\le\bigl(D_hA(h)e_1,e_1\bigr)=\bigl(D_hA(h)e_1\bigr)^1
\le\bigl|\bigl(D_hA(h)e_1\bigr)^1\bigr|
\le\gamma_1(\varepsilon+1)^{\frac p2-1},
\]
and $(\varepsilon+1)^{p/2-1}>0$, so $\gamma_0\le\gamma_1$ and hence
$0<\gamma_1$.

\emph{Claim A (the nondegenerate case).}  \emph{Let $\varepsilon>0$,
$\gamma_0>0$ and $\gamma_1\le\bar r\gamma_0$ with $\gamma_1>0$, let
$\Omega'\subseteq\R^n$ be open, let $A'\colon\R^n\to\R^n$ satisfy
(Q1) of the thesis--(Q2) of the thesis on $\Omega'$ with parameters
$p,\varepsilon,\gamma_0,\gamma_1$ (\noderef{EllipticStructureQ12}), let $u'$ be a
weak solution of $\operatorname{div}A'(\nabla u')=0$ in $\Omega'$
(\noderef{IsWeakSolutionOn}) with weak gradient $g'$, and let $y\in\R^n$ and
$\rho>0$ satisfy $\overline B_\rho(y)\subseteq\Omega'$.  Then for almost every
$x\in B_{\rho/2}(y)$,}
\[
\bigl(\varepsilon+|g'(x)|^2\bigr)^{\frac p2}
\ \le\ \frac{C_1}{\rho^n}\int_{B_\rho(y)}\bigl(\varepsilon+|g'|^2\bigr)^{\frac p2}dx .
\tag{A}
\]

\emph{Proof of Claim A, step A1: what (C1) supplies.}  The present theorem
assumes the cited regularity result (C1) for every field, every open set and all
parameters with positive structure constants.  Instantiate that clause at
$\varepsilon,\gamma_0,\gamma_1,\Omega',A'$: its side conditions
$\gamma_0>0$, $\gamma_1>0$ and openness of $\Omega'$ hold.  The resulting
predicate \noderef{CitedRegularityEpsPos} has as its own antecedents
$\varepsilon>0$, the structure conditions on $\Omega'$ and the weak-solution
property, all of which hold, and it yields
\begin{enumerate}
\item $u'\in W^{1,2}_{\mathrm{loc}}(\Omega')$ with weak gradient $g'$
  (\noderef{MemW1pLocOn}); and
\item a field $\hat g$ with (i) $\hat g$ continuous on $\Omega'$, (ii)
  $g'=\hat g$ almost everywhere on $\Omega'$, and (iii) a family
  $D\hat g=(D\hat g^1,\dots,D\hat g^n)$ with
  $\hat g^k\in W^{1,2}_{\mathrm{loc}}(\Omega')$ and weak gradient $D\hat g^k$
  for each $k$ (\noderef{MemW1pLocOn}).
\end{enumerate}
Fix such a $\hat g$ and such a family $D\hat g$ once and for all.

\emph{A warning about the prime, because the notation collides.}  In this proof
$g'$ is \emph{Claim A's own given field} --- the weak gradient of $u'$, named in
the claim's hypotheses above --- and it is \textbf{not} the representative.
\noderef{CitedRegularityEpsPos} happens to write $g'$ for \emph{its}
representative, so its clauses (1)--(3) are transcribed here with $\hat g$ in
that role, to keep the two apart.  Clauses (i) and (iii) are properties of
$\hat g$ and \emph{not} of $g'$: a weak gradient is an almost-everywhere class,
so neither continuity nor pointwise boundedness can be asserted of $g'$, and
that node's own commentary records that a version asserting continuity of the
given field would be a false $\mathrm{Prop}$.  \emph{The consequence for this
proof is structural, not cosmetic: $w$ and $v$ below are built from $\hat g$,
steps A2--A6 are run for $\hat g$, and the passage back to $g'$ is made once, at
A7, where the conclusion is already an almost-everywhere statement.}  This is
available here precisely because $w$ and $v$ are internal to Claim A --- no
later step or cited node fixes their names --- so the representative may be
carried where points are needed and the given field where only integrals are.

\emph{A2: the functions $w$ and $v$, built from the representative.}  Put
\[
w=\varepsilon+|\hat g|^2,\qquad v=w^{\frac p4},
\]
both defined at every point.  On $\Omega'$ we have $w\ge\varepsilon>0$, so
$v>0$ there; and $w$ is continuous on $\Omega'$ by A1(2)(i), hence so is $v$,
because $s\mapsto s^{p/4}$ is continuous on $(0,\infty)$.  Note for later that
$v^2=w^{p/2}$ and, for $\theta>0$, $(v^\theta)^2=w^{p\theta/2}$, the exponent
rules for real powers being available since $w>0$ on $\Omega'$.

\emph{A3: $v^\theta$ is a local $W^{1,2}$ function, with an explicit weak
gradient.}  Fix $t>0$ and put $\theta=1+\frac tp$, so that $\theta>1$ and
$v^\theta=w^{\frac{p+t}4}$; as $t$ runs over $(0,\infty)$, $\theta$ runs over all
of $(1,\infty)$, and $t=p(\theta-1)$ recovers $t$ from $\theta$.  Apply
\noderef{WeakGradientVectorComp} with $n$, $m=n$, $q=2$, the open set $\Omega'$,
the functions $G_k=\hat g^k$ with weak gradients $DG_k=D\hat g^k$, and
\[
F(y_1,\dots,y_n)=\Bigl(\varepsilon+\sum_{k=1}^ny_k^2\Bigr)^{\frac{p+t}4}.
\]
Its hypotheses hold: $1\le n$ and $1\le 2$; $\Omega'$ is open; each
$G_k\in W^{1,2}_{\mathrm{loc}}(\Omega')$ with weak gradient $DG_k$ by
A1(2)(iii); the family is locally uniformly bounded, since for compact
$K\subseteq\Omega'$ the function $|\hat g|$ is continuous on $K$ by A1(2)(i) and
therefore bounded there, say by $L$, whence $|G_k(x)|\le|\hat g(x)|\le L$ for
all $x\in K$ and all $k$.  \emph{This is the hypothesis that forces the
representative on us: it asks for a bound at \emph{every} point of $K$, and the
given $g'$ need not satisfy one, a weak gradient being alterable on a null set
without disturbing any hypothesis of this claim.}  And $F$ is of
class $C^1$ on $\R^n$, being the composition of the polynomial
$y\mapsto\varepsilon+|y|^2$, which is smooth and takes values in
$[\varepsilon,\infty)\subseteq(0,\infty)$ because $\varepsilon>0$, with
$s\mapsto s^{(p+t)/4}$, which is smooth on $(0,\infty)$.  Since
$F(G(x))=\bigl(\varepsilon+|\hat g(x)|^2\bigr)^{(p+t)/4}=v^\theta(x)$, the conclusion
is that $v^\theta\in W^{1,2}_{\mathrm{loc}}(\Omega')$ with weak gradient
\[
G_\theta=\sum_{k=1}^nD_kF(G)\,D\hat g^k
=\frac{p+t}4\,w^{\frac{p+t}4-1}\,\nabla w,
\qquad
\nabla w:=2\sum_{k=1}^n\hat g^k\,D\hat g^k ,
\]
where the middle equality is the computation
$D_kF(y)=\frac{p+t}4\bigl(\varepsilon+|y|^2\bigr)^{\frac{p+t}4-1}\cdot2y_k$
evaluated at $y=G(x)$.  Finally $B_\rho(y)$ is open and contained in
$\overline B_\rho(y)\subseteq\Omega'$, so clause (2) of
\noderef{WeakGradientRestrict}, applied with $q=2$, the sets
$\Omega'\supseteq B_\rho(y)$ and the pair $(v^\theta,G_\theta)$, gives
\[
v^\theta\in W^{1,2}_{\mathrm{loc}}\bigl(B_\rho(y)\bigr)
\quad\text{with weak gradient}\quad G_\theta .
\tag{A3}
\]

\emph{A4: the Caccioppoli inequality in the form Moser's lemma consumes.}
Apply \noderef{CaccioppoliPowerWeight} for the triple $(n,p,\bar r)$, whose
constant $c$ was fixed in Step 0, to the data
$\varepsilon,\gamma_0,\gamma_1,t,\Omega',A',u',g',D\hat g$ and an arbitrary
$\zeta$ which is of class $C^\infty$ on $\R^n$, has compact support, and
satisfies $\mathrm{tsupport}\,\zeta\subseteq\Omega'$ --- these are the three
conditions that node binds on its cut-off, and $\mathrm{tsupport}$, the
\emph{closure} of the support, is the one it asks about, so it is the one used
here and below.  Its hypotheses are $\varepsilon>0$,
$\gamma_0>0$, $\gamma_1\le\bar r\gamma_0$, $t>0$, openness of $\Omega'$, the
structure conditions (\noderef{EllipticStructureQ12}), the cited regularity
result (C1) (\noderef{CitedRegularityEpsPos}) at these very data, which is the
instance A1 already formed, the weak-solution property
(\noderef{IsWeakSolutionOn}), the membership
$g'^k\in W^{1,2}_{\mathrm{loc}}(\Omega')$ with weak gradient $D\hat g^k$ for
every $k$ (\noderef{MemW1pLocOn}), established in the next paragraph, and the
three cut-off conditions
just listed, namely $\mathrm{ContDiff}\;\R\;\infty\;\zeta$,
$\mathrm{HasCompactSupport}\;\zeta$ and
$\mathrm{tsupport}\,\zeta\subseteq\Omega'$.

\emph{The one transfer back to the given field.}  That node carries its
gradient-regularity hypothesis on the \emph{same} field as its weak-solution
hypothesis, namely on $g'$, whereas A1(2)(iii) supplies it for $\hat g$.  The
two are reconciled once, and the fact needed is that
$\mathrm{MemW1pLocOn}$ is unchanged when the function is modified on a null set,
the weak gradient being kept.  By \noderef{MemW1pLocOn} that assertion is two
things, and each is an integral one.  First the identity of
\noderef{HasWeakGradientOn}: for every test function $\psi$ of that node's class
and every index $i$,
$\int_{\Omega'}\hat g^k\,D_i\psi\,dx=-\int_{\Omega'}(D\hat g^k)_i\,\psi\,dx$,
and $g'^k=\hat g^k$ almost everywhere on $\Omega'$ by A1(2)(ii), so the two
left-hand integrands agree almost everywhere and the integrals are equal,
giving $\int_{\Omega'}g'^k\,D_i\psi\,dx=-\int_{\Omega'}(D\hat g^k)_i\,\psi\,dx$.
Second the $L^2$ clauses: for compact $K\subseteq\Omega'$ we have
$g'^k=\hat g^k$ almost everywhere on $K$ and $D\hat g^k$ is not changed at all,
so $g'^k\in L^2(K)$ and $D\hat g^k\in L^2(K)$ follow from the same statements
for $\hat g^k$.  Hence
\[
g'^k\in W^{1,2}_{\mathrm{loc}}(\Omega')\ \text{ with weak gradient }\ D\hat g^k
\qquad(k=1,\dots,n),
\tag{A4a}
\]
which is the hypothesis just invoked.  (No further node is needed for it: the
two clauses of \noderef{MemW1pLocOn} are the only facts used, and both are
integral assertions.)

Its conclusion is
\[
\int_{\Omega'}\Bigl(\frac{p+t}4\Bigr)^2w_\flat^{\frac{p+t}2-2}
\Bigl(\sum_{i=1}^n(\nabla w_\flat)_i^2\Bigr)\zeta^2\,dx
\ \le\ c\int_{\Omega'}w_\flat^{\frac{p+t}2}\,|\nabla\zeta|^2\,dx ,
\]
where $w_\flat=\varepsilon+|g'|^2$ and
$(\nabla w_\flat)_i=2\sum_kg'^k(D\hat g^k)^i$ are the weight and the field that
node forms from the data it was given, i.e.\ from $g'$ and $D\hat g$.

\emph{The weights agree almost everywhere, so the display may be rewritten in
$w$.}  By A1(2)(ii), $g'=\hat g$ almost everywhere on $\Omega'$, and
$w_\flat$, $\nabla w_\flat$ are obtained from $g'$ by the same two formulas that
produce $w$, $\nabla w$ from $\hat g$ --- the second with the \emph{same} family
$D\hat g$ in both cases --- so
\[
w_\flat=w\quad\text{and}\quad \nabla w_\flat=\nabla w
\qquad\text{almost everywhere on }\Omega' ,
\tag{A4b}
\]
and hence so does every real power of the weight.  Both sides of the display are
integrals, so altering their integrands on a null set changes neither, and the
display holds verbatim with $w$ and $\nabla w$ in place of $w_\flat$ and
$\nabla w_\flat$.  Now by A3
\[
|G_\theta|^2=\Bigl(\frac{p+t}4\Bigr)^2\bigl(w^{\frac{p+t}4-1}\bigr)^2|\nabla w|^2
=\Bigl(\frac{p+t}4\Bigr)^2w^{\frac{p+t}2-2}\sum_{i=1}^n(\nabla w)_i^2
\]
on $\Omega'$, using $(w^a)^2=w^{2a}$ for $w>0$ and
$2\bigl(\frac{p+t}4-1\bigr)=\frac{p+t}2-2$; and
$(v^\theta)^2=w^{p\theta/2}=w^{\frac{p+t}2}$ by A2.  So the display reads
\[
\int_{\Omega'}|G_\theta|^2\zeta^2\,dx
\ \le\ c\int_{\Omega'}(v^\theta)^2|\nabla\zeta|^2\,dx
\]
for every such $\zeta$.
If moreover $\mathrm{tsupport}\,\zeta\subseteq B_\rho(y)$, then both integrands
vanish at every point outside $\mathrm{tsupport}\,\zeta$: there $\zeta=0$,
because $\mathrm{supp}\,\zeta\subseteq\mathrm{tsupport}\,\zeta$, so $\zeta^2=0$;
and the complement of $\mathrm{tsupport}\,\zeta$ is open with $\zeta$
identically $0$ on it, so $\nabla\zeta=0$ there as well.  Hence each integrand
has the same product with the indicator of $\Omega'$ as with the indicator of
$B_\rho(y)$, the two indicators agreeing on $\mathrm{tsupport}\,\zeta$, and the
two integrals may be taken over $B_\rho(y)$:
\[
\int_{B_\rho(y)}|G_\theta|^2\zeta^2\,dx
\ \le\ c\int_{B_\rho(y)}(v^\theta)^2|\nabla\zeta|^2\,dx
\tag{A4}
\]
for every $\zeta$ of class $C^\infty$ with compact support and
$\mathrm{tsupport}\,\zeta\subseteq B_\rho(y)$ --- which is exactly the cut-off
class the Caccioppoli hypothesis of \noderef{MoserLocalBoundedness} quantifies
over, at the ball $B_\rho(y)$.  The constant
$c$ does not depend on $t$, hence not on $\theta$: it was chosen in Step 0,
before $t$ was introduced, and this uniformity is exactly what
\noderef{MoserLocalBoundedness} requires of its $Q$.  A $\theta$-dependent
constant would not do, and would also not leave $C$ on the dependency list the
statement prescribes.

\emph{A5: truncation, and the three remaining hypotheses of Moser's lemma.}  Put
$\tilde v=\mathbf1_{B_\rho(y)}v$, the function equal to $v$ on $B_\rho(y)$ and to
$0$ elsewhere.  Then:
\begin{enumerate}
\item $\tilde v\ge0$ everywhere, since $v=w^{p/4}\ge0$ as a real power of a
  nonnegative number.
\item $\tilde v$ is measurable.  For $a<0$ the superlevel set
  $\{\tilde v>a\}$ is all of $\R^n$; for $a\ge0$ it is
  $\{x\in B_\rho(y):v(x)>a\}$, which is open because $v$ is continuous on the
  open set $B_\rho(y)\subseteq\Omega'$ by A2 --- legitimately, $v$ having been
  built from the continuous $\hat g$ and not from $g'$.  All superlevel sets being
  measurable, $\tilde v$ is measurable (Preamble).
\item $\tilde v^2$ is integrable on $B_\rho(y)$.  Indeed
  $\tilde v^2=w^{p/2}$ there, $\overline B_\rho(y)$ is a compact subset of
  $\Omega'$ and $w$ is continuous on $\Omega'$, so $w^{p/2}$ is continuous on
  $\overline B_\rho(y)$ and therefore bounded there, while
  $|B_\rho(y)|<\infty$; the Preamble item on domination gives integrability.
\end{enumerate}
Moreover $\tilde v^\theta=v^\theta$ on $B_\rho(y)$ for every $\theta>0$.  The two
hypotheses of \noderef{MoserLocalBoundedness} that mention the powers of the
function are statements about $B_\rho(y)$ only: the defining identity of
\noderef{HasWeakGradientOn} on $B_\rho(y)$ and the local $L^2$ clauses of
\noderef{MemW1pLocOn} on $B_\rho(y)$ integrate over $B_\rho(y)$ and over its
compact subsets, and (A4) integrates over $B_\rho(y)$.  Since $\tilde v^\theta$
and $v^\theta$ agree there, (A3) and (A4) hold verbatim with $\tilde v$ in place
of $v$.

\emph{A6: Moser's lemma.}  Apply \noderef{MoserLocalBoundedness} with $n$,
$Q=c$, the hypothesis \noderef{CitedSobolevPoincareBall} of the present theorem,
and then with the centre $y$, the radius $\rho$, the function $\tilde v$ and the
family of fields $\theta\mapsto G_{p(\theta-1)}$ for $\theta>1$ (and, say, $0$
for $\theta\le1$; the hypotheses only ever evaluate it at $\theta>1$).  Its
hypotheses are: $2\le n$ and $0<c$ (Step 0); (C6), which the present theorem
assumes; $0<\rho$; measurability and nonnegativity of $\tilde v$ and
integrability of $\tilde v^2$ on $B_\rho(y)$, which are A5(1)--(3); the
membership $\tilde v^\theta\in W^{1,2}_{\mathrm{loc}}(B_\rho(y))$ with weak
gradient $G_{p(\theta-1)}$ for every $\theta>1$, which is (A3) together with
A5, since $t=p(\theta-1)>0$ and then $\theta=1+t/p$; and the Caccioppoli
inequality for every $\theta>1$ and every cut-off, with the single constant $c$,
which is (A4) together with A5.  The conclusion is: for almost every
$x\in B_{\rho/2}(y)$,
\[
\tilde v(x)\ \le\ C_0\Bigl(\rho^{-n}\int_{B_\rho(y)}\tilde v^2\,dx\Bigr)^{\frac12}.
\]

\emph{A7: squaring, and the passage back to the given field.}  Let
$x\in B_{\rho/2}(y)$ be a point where the last display holds \emph{and} at which
$g'(x)=\hat g(x)$; almost every point of $B_{\rho/2}(y)$ is such a point, being
in the intersection of two sets of full measure in $B_\rho(y)$ --- the first
from A6's conclusion, the second from A1(2)(ii).  There
$\tilde v(x)=v(x)=w(x)^{p/4}$, since $B_{\rho/2}(y)\subseteq B_\rho(y)$, and
$w(x)=\varepsilon+|\hat g(x)|^2=\varepsilon+|g'(x)|^2$ by the choice of $x$.
Both sides of the display are nonnegative, so squaring preserves the inequality,
and $\bigl(w(x)^{p/4}\bigr)^2=w(x)^{p/2}$ while
$\int_{B_\rho(y)}\tilde v^2=\int_{B_\rho(y)}w^{p/2}$ by A5(3); and
$\int_{B_\rho(y)}w^{p/2}=\int_{B_\rho(y)}(\varepsilon+|g'|^2)^{p/2}$, the two
integrands agreeing almost everywhere by (A4b).  \emph{This is the single point
at which the argument returns to $g'$, and it is legitimate because the
conclusion (A) is itself an almost-everywhere statement and its right-hand side
is an integral.}  Hence for almost
every $x\in B_{\rho/2}(y)$
\[
\bigl(\varepsilon+|g'(x)|^2\bigr)^{\frac p2}=w(x)^{\frac p2}
\ \le\ \frac{C_0^2}{\rho^n}\int_{B_\rho(y)}w^{\frac p2}\,dx
=\frac{C_1}{\rho^n}\int_{B_\rho(y)}\bigl(\varepsilon+|g'|^2\bigr)^{\frac p2}dx ,
\]
which is (A).  This proves Claim A.

\emph{Step 2: the theorem when $\varepsilon>0$.}  Apply Claim A with
$\Omega'=\Omega$, $A'=A$, $u'=u$, $g'=g$, $y=x_0$ and $\rho=R$.  Its hypotheses
hold: $\varepsilon>0$; $\gamma_0>0$; $\gamma_1\le r\gamma_0\le\bar r\gamma_0$,
because $r\le\bar r$ and $\gamma_0>0$; $\gamma_1>0$ by Step 1; $\Omega$ open;
the structure conditions and the weak-solution property, which are hypotheses of
the theorem; $R>0$; and $\overline B_R\subseteq\Omega$, which is the containment
hypothesis.  So for almost every $x\in B_{R/2}$
\[
\bigl(\varepsilon+|g(x)|^2\bigr)^{\frac p2}
\le\frac{C_1}{R^n}\int_{B_R}\bigl(\varepsilon+|g|^2\bigr)^{\frac p2}dx
\le\frac{C}{R^n}\int_{B_R}\bigl(\varepsilon+|g|^2\bigr)^{\frac p2}dx ,
\]
the last step because $C_1\le C$ and the integral is nonnegative, its integrand
being a real power of a nonnegative number.  This is the assertion.

\emph{Step 3: the theorem when $\varepsilon=0$.}  Assume now $\varepsilon=0$.

\emph{B1: the structure conditions hold on all of $\R^n$.}  The field of the
present theorem is $(x,z,h)\mapsto A(h)$, independent of $x$ and $z$.
Two of the six clauses of \noderef{EllipticStructureQ12} make no reference to
the domain at all --- Carath\'eodory measurability in the point asserts that for
every $z$ and $h$ the map $x\mapsto A(x,z,h)$ is Borel measurable on $\R^n$, and
Carath\'eodory joint continuity asserts that for every $x\in\R^n$ the map
$(z,h)\mapsto A(x,z,h)$ is continuous --- so each is literally the same
assertion whichever domain the predicate is read with, and both carry over
verbatim rather than by any argument about points.  The remaining four are the ones that have to be moved, and read
for this field, those four clauses of \noderef{EllipticStructureQ12} on $\Omega$ say:
$A$ is continuous on $\R^n$ and of class $C^1$ on
$D(0)=\{h:|h|^2\ne0\}$, for every $x\in\Omega$ and every $z$; and
(Q1) of the thesis--(Q2) of the thesis hold, for almost every $x\in\Omega$, for every $z$,
every $h\in D(0)$ and every $\xi$.  None of the four assertions mentions $x$ or
$z$ once the field is written out.  The set $\Omega$ contains $\overline B_R$
with $R>0$, so it is nonempty and of positive measure (Preamble: the volume of
a ball); the two almost-everywhere clauses each fail only on a null set, the
union of those two null sets is null, and a null set cannot contain the
positive-measure $\Omega$, so there is a point $x\in\Omega$ at which both hold;
and the first two clauses hold at every point of $\Omega$, being asserted
there.  Therefore all four of those assertions hold
outright, and with the two domain-free clauses they are all six, and in particular they hold for every $x\in\R^n$ and almost every
$x\in\R^n$ respectively, i.e.
\[
\text{$A$ satisfies (Q1) of the thesis--(Q2) of the thesis on $\R^n$ with parameters
$p,0,\gamma_0,\gamma_1$ (\noderef{EllipticStructureQ12}).}
\tag{B1}
\]

\emph{B2: $u$ is a weak solution in the ball $B_R$.}  By
\noderef{IsWeakSolutionOn} this has three parts: the membership clause, and,
for each admissible test pair at the ball, the integrability of the integrand on
the ball together with the vanishing of its integral.  First,
$u\in W^{1,p}_{\mathrm{loc}}(B_R)$ with weak gradient $g$: this is clause (2) of
\noderef{WeakGradientRestrict} applied with $q=p$, the sets
$\Omega\supseteq B_R$ --- open, with $B_R\subseteq\overline B_R\subseteq\Omega$
--- and the pair $(u,g)$, whose hypothesis
$u\in W^{1,p}_{\mathrm{loc}}(\Omega)$ with weak gradient $g$ is the first clause
of \noderef{IsWeakSolutionOn} on $\Omega$.  Second, the integral identity for
the test class of $B_R$: let $\varphi\in W^{1,p}(B_R)$ with weak gradient
$g_\varphi$ (\noderef{MemW1pOn}) have compact support with
$\mathrm{tsupport}\,\varphi\subseteq B_R$, these being the three conditions the
test clause of \noderef{IsWeakSolutionOn} binds at the ball.  Apply
\noderef{WeakGradientZeroExtend} with $q=p$, the open set $\Omega$, the ball
$B_R\subseteq\Omega$ and the pair $(\varphi,g_\varphi)$: it gives
$\varphi\in W^{1,p}(\Omega)$ with weak gradient
$G_\varphi=\mathbf1_{B_R}g_\varphi$ --- its own hypotheses being $1\le n$,
$1\le p$, the openness of $\Omega$, the inclusion $B_R\subseteq\Omega$, the
membership just assumed, the compactness of $\mathrm{supp}\,\varphi$ and
$\mathrm{tsupport}\,\varphi\subseteq B_R$ --- and $\varphi$ has compact support
with $\mathrm{tsupport}\,\varphi\subseteq B_R\subseteq\Omega$, which is the
inclusion the test clause on $\Omega$ requires.  So
$(\varphi,G_\varphi)$ belongs to the test class of \noderef{IsWeakSolutionOn} on
$\Omega$, and the identity there reads
\[
\int_\Omega\Bigl[\sum_{i=1}^nA^i\bigl(g(x)\bigr)G_\varphi^i(x)+0\cdot\varphi(x)\Bigr]dx=0 .
\]
The integrand vanishes at every point outside $B_R$: there $G_\varphi=0$,
being the zero extension of $g_\varphi$, and $\varphi=0$, because
$\mathrm{supp}\,\varphi\subseteq\mathrm{tsupport}\,\varphi\subseteq B_R$.  So the
integral over $\Omega$ equals the integral over $B_R$, where
$G_\varphi=g_\varphi$.  Hence
$\int_{B_R}\bigl[\sum_iA^i(g)g_\varphi^i+0\cdot\varphi\bigr]dx=0$.

Third, the integrability clause at the ball.  The test clause of
\noderef{IsWeakSolutionOn} on $\Omega$, applied to the pair
$(\varphi,G_\varphi)$ --- admissible there by the previous paragraph --- asserts
not only the vanishing of the integral but also that the $\Omega$-integrand
$x\mapsto\sum_iA^i(g(x))G_\varphi^i(x)+0\cdot\varphi(x)$ is integrable on
$\Omega$.  The ball $B_R$ is a measurable subset of $\Omega$, and the
restriction of an integrable function to a measurable subset is integrable, so
that integrand is integrable on $B_R$; and on $B_R$ it coincides with the
$B_R$-integrand $x\mapsto\sum_iA^i(g(x))g_\varphi^i(x)+0\cdot\varphi(x)$,
because $G_\varphi=\mathbf1_{B_R}g_\varphi$ equals $g_\varphi$ there.  Hence the
$B_R$-integrand is integrable on $B_R$, which is the clause required.

As the test pair was arbitrary, $u$ is a weak solution of
$\operatorname{div}A(\nabla u)=0$ in $B_R$ (\noderef{IsWeakSolutionOn}).

\emph{B3: the approximating sequence.}  Apply
\noderef{DegenerateLimitTransfer} with $n,p,r$ --- the same instance as in
Step 0, so the number it supplies is the $r'$ of Step 0 --- and then with
$\gamma_0,\gamma_1$ ($0<\gamma_0$ and $\gamma_1\le r\gamma_0$ are hypotheses of
the theorem), with the field $A$, for which $A(0)=0$ is a hypothesis of the
theorem and (B1) is the structure requirement, with the cited solvability result
(C2) for every field and all parameters with positive structure constants, which
is a hypothesis of the present theorem in exactly that form, and with
$x_0,R,u,g$, for which $R>0$ holds, B2 gives the weak-solution requirement, and
its remaining hypothesis $u\in W^{1,p}(B_R)$ with weak gradient $g$
(\noderef{MemW1pOn}) is clause (4) of \noderef{WeakGradientRestrict} applied
with $q=p$, the sets $\Omega$ and $\Omega'=B_R$ and the pair $(u,g)$: its
hypotheses there are $1\le n$, $1\le p$, openness of $\Omega$ and of $B_R$,
$B_R\subseteq\Omega$, compactness of $\overline{B_R}$ --- which is the closure of
$B_R$ because $R>0$ --- together with $\overline{B_R}\subseteq\Omega$, the
containment hypothesis of the theorem, and
$u\in W^{1,p}_{\mathrm{loc}}(\Omega)$ with weak gradient $g$, the first clause of
\noderef{IsWeakSolutionOn} on $\Omega$.  It
supplies a sequence $e_m$, fields $A_{e_m}$ and functions $u_m$ with weak
gradients $g_m$ such that
\begin{enumerate}
\item $e_m>0$ for every $m$ and $e_m\to0$;
\item for every $m$, $A_{e_m}(0)=0$ and there are $\gamma_0^{(m)}>0$ and
  $\gamma_1^{(m)}\le r'\gamma_0^{(m)}$ for which $A_{e_m}$ satisfies
  (Q1) of the thesis--(Q2) of the thesis on $\R^n$ with parameter $e_m$ and those
  constants (\noderef{EllipticStructureQ12});
\item $u_m$ is a weak solution of $\operatorname{div}A_{e_m}(\nabla u_m)=0$ in
  $B_R$ with weak gradient $g_m$ (\noderef{IsWeakSolutionOn});
\item $\int_{B_R}|g_m-g|^p\,dx\to0$; and
\item $g_m\to g$ almost everywhere on $B_R$.
\end{enumerate}
(The clause of that node about $u_m-u\in W^{1,p}_0(B_R)$ is not used here.)

\emph{B4: the estimate for each $m$ on the smaller balls.}  Fix
$\rho\in(0,R)$ and $m\in\mathbb{N}$.  Apply Claim A with $\varepsilon:=e_m>0$,
$\gamma_0:=\gamma_0^{(m)}$, $\gamma_1:=\gamma_1^{(m)}$, $\Omega':=B_R$,
$A':=A_{e_m}$, $u':=u_m$, $g':=g_m$, $y:=x_0$ and the radius $\rho$.  Its
hypotheses hold: $e_m>0$ by B3(1); $\gamma_0^{(m)}>0$ by B3(2);
$\gamma_1^{(m)}\le r'\gamma_0^{(m)}\le\bar r\gamma_0^{(m)}$ since
$r'\le\bar r$ and $\gamma_0^{(m)}>0$; $\gamma_1^{(m)}>0$, by the argument of
Step 1 applied to the field $A_{e_m}$ on the nonempty open set $B_R$ --- which is
legitimate because B3(2) gives (Q1) of the thesis--(Q2) of the thesis for $A_{e_m}$ on
$\R^n$ and therefore on $B_R$, the six clauses of
\noderef{EllipticStructureQ12} being inherited by subsets: the first two, the
Carath\'eodory pair, mention no domain and so are unchanged, and the remaining
four are each an assertion about all, respectively almost all, points of the
set, hence weaken when the set shrinks; $B_R$ is open;
the structure conditions on $B_R$, as just noted; the weak-solution property on
$B_R$, which is B3(3); $\rho>0$; and
$\overline B_\rho(x_0)\subseteq B_R$, because $\rho<R$.  So for almost every
$x\in B_{\rho/2}$
\[
\bigl(e_m+|g_m(x)|^2\bigr)^{\frac p2}
\ \le\ \frac{C_1}{\rho^n}\int_{B_\rho}\bigl(e_m+|g_m|^2\bigr)^{\frac p2}dx .
\tag{B4}
\]

\emph{B5: the limit in $m$.}  We bound the right-hand side of (B4) above, up to a
factor, by the integral of $|g|^p$.  For $a,b\ge0$ one has
\[
(a+b)^{\frac p2}\le\kappa\bigl(a^{\frac p2}+b^{\frac p2}\bigr),
\qquad \kappa=\max\{1,2^{\frac p2-1}\} :
\]
for $p\le2$ the exponent $p/2$ lies in $(0,1]$ and $s\mapsto s^{p/2}$ is
subadditive on $[0,\infty)$, so the inequality holds with $1$; for $p\ge2$ it is
the convexity of $s\mapsto s^{p/2}$ in the form
$\bigl(\frac{a+b}2\bigr)^{p/2}\le\frac12\bigl(a^{p/2}+b^{p/2}\bigr)$.  Applying
this twice, together with $|g_m|^2\le2|g|^2+2|g_m-g|^2$ --- which is
$|g_m|\le|g|+|g_m-g|$ followed by $(s+u)^2\le2s^2+2u^2$ --- gives pointwise
on $B_\rho$
\[
\bigl(e_m+|g_m|^2\bigr)^{\frac p2}
\le\kappa\Bigl(e_m^{\frac p2}+\bigl(2|g|^2+2|g_m-g|^2\bigr)^{\frac p2}\Bigr)
\le\kappa e_m^{\frac p2}+\kappa^2 2^{\frac p2}\bigl(|g|^p+|g_m-g|^p\bigr),
\]
where in the last step $\bigl(2(|g|^2+|g_m-g|^2)\bigr)^{p/2}
=2^{p/2}\bigl(|g|^2+|g_m-g|^2\bigr)^{p/2}
\le2^{p/2}\kappa\bigl(|g|^p+|g_m-g|^p\bigr)$, using
$\bigl(|g|^2\bigr)^{p/2}=|g|^p$.  The two displayed inequalities hold at every
point of $B_R$, not merely on $B_\rho$.  All the functions appearing are
nonnegative and measurable, and $|g|^p$ and $|g_m-g|^p$ are integrable on $B_R$ --- the former
because $\overline B_R\subseteq\Omega$ is compact and the first clause of
\noderef{IsWeakSolutionOn} puts $|g|$ in $L^p$ of the compact subsets of
$\Omega$ (\noderef{MemW1pLocOn}), whence $|g|\in L^p(B_R)$ because
$|g|^p\mathbf 1_{B_R}\le|g|^p\mathbf 1_{\overline B_R}$ pointwise with the
right-hand side integrable, so clause (a) of the Preamble's domination
proposition applies; the latter from the clause of
\noderef{DegenerateLimitTransfer} about $u_m-u\in W^{1,p}_0(B_R)$, which B3 set
aside and which is needed here.  By \noderef{MemW1pZeroOn} that clause contains,
as its first conjunct, $u_m-u\in W^{1,p}(B_R)$ with weak gradient $g_m-g$
(\noderef{MemW1pOn}), and the $L^p$ clause of \noderef{MemW1pOn} is exactly
$|g_m-g|\in L^p(B_R)$, i.e.\ $|g_m-g|^p$ integrable on $B_R$.  \emph{B3(4) will
not serve in its place: it says only that $\int_{B_R}|g_m-g|^p\,dx\to0$, and the
integral of a non-integrable function is $0$, so a vanishing or finite integral
does not by itself give integrability --- which is the very reason
\noderef{IsWeakSolutionOn} carries an explicit integrability clause.}  Now for a nonnegative
measurable $h$ and measurable $E\subseteq F$ one has $h\mathbf 1_E\le h\mathbf
1_F$ pointwise, so $\int_Eh\,dx\le\int_Fh\,dx$ whenever $h$ is integrable on $F$
(this is monotonicity of the integral in the integrand on the fixed set $F$;
Mathlib: \texttt{MeasureTheory.setIntegral\_mono\_set}).  Write
$W=\kappa e_m^{p/2}+\kappa^22^{p/2}\bigl(|g|^p+|g_m-g|^p\bigr)$ for the
right-hand side of the pointwise bound; $W$ is integrable on $B_R$, being a
constant --- integrable because $|B_R|<\infty$ (Preamble: the volume of a
ball) --- plus the two terms just shown
integrable, so by clause (a) again $\bigl(e_m+|g_m|^2\bigr)^{p/2}$ is integrable
on $B_R$ as well.  Taking $h=\bigl(e_m+|g_m|^2\bigr)^{p/2}$, $E=B_\rho$ and
$F=B_R$ in the previous sentence, and then bounding $\int_{B_R}h\,dx$ by
$\int_{B_R}W\,dx$ by monotonicity of the integral in the integrand on the fixed
set $B_R$ --- both functions being integrable there and the inequality holding
pointwise (Mathlib: \texttt{MeasureTheory.setIntegral\_mono\_on}) --- gives
\[
\int_{B_\rho}\bigl(e_m+|g_m|^2\bigr)^{\frac p2}dx
\ \le\ \kappa\,e_m^{\frac p2}\,|B_R|
+\kappa^2 2^{\frac p2}\Bigl(\int_{B_R}|g|^p\,dx+\int_{B_R}|g_m-g|^p\,dx\Bigr).
\]
As $m\to\infty$ the first term tends to $0$: $e_m\to0$ by B3(1), and the real
power $s\mapsto s^{p/2}$ is continuous on $[0,\infty)$ with $0^{p/2}=0$, the
exponent $p/2$ being positive and nonzero (the conventions for real powers are
those recorded in the Preamble).  The last term tends to $0$ by B3(4).  Hence
\[
\limsup_{m\to\infty}\int_{B_\rho}\bigl(e_m+|g_m|^2\bigr)^{\frac p2}dx
\ \le\ \kappa^2 2^{\frac p2}\int_{B_R}|g|^p\,dx .
\tag{B5}
\]
On the left-hand side of (B4), let $x\in B_{\rho/2}$ be a point at which
$g_m(x)\to g(x)$; by B3(5) almost every point of $B_R$, hence almost every point
of $B_{\rho/2}$, is such a point.  Then
$e_m+|g_m(x)|^2\to0+|g(x)|^2$, and since the real power $s\mapsto s^{p/2}$ is continuous on
$[0,\infty)$ --- again because $p/2>0$, which is what makes it continuous at $0$
as well, with $0^{p/2}=0$ --- we get
\[
\bigl(e_m+|g_m(x)|^2\bigr)^{\frac p2}\longrightarrow
\bigl(0+|g(x)|^2\bigr)^{\frac p2}=|g(x)|^p .
\]
Now fix $\rho\in(0,R)$ and collect the exceptional sets.  (B4) was obtained
from Claim A \emph{one $m$ at a time}, so for each $m$ it holds at every
$x\in B_{\rho/2}$ outside a null set $Z_m$; and the convergence
$g_m(x)\to g(x)$ holds outside one further null set $Z_\infty$, by B3(5).  Put
$Z=Z_\infty\cup\bigcup_{m\in\mathbb{N}}Z_m$, a null set, being a countable union of null
sets (Preamble: the measurable sets are closed under countable unions, and a
countable union of null sets is null by countable subadditivity of the measure).
For every $x\in B_{\rho/2}\setminus Z$ both (B4) holds for every $m$ and the
left-hand side converges to $|g(x)|^p$; and almost every point of $B_{\rho/2}$
is such a point.  Passing to the limit in (B4) along such an $x$ and
using (B5),
\[
\bigl(0+|g(x)|^2\bigr)^{\frac p2}\ \le\ \frac{C_1}{\rho^n}\,\kappa^2 2^{\frac p2}
\int_{B_R}|g|^p\,dx
=\frac{C}{\rho^n}\int_{B_R}\bigl(0+|g|^2\bigr)^{\frac p2}dx
\tag{B6}
\]
for almost every $x\in B_{\rho/2}$, where $C=C_1\kappa^22^{p/2}$ is the constant
of Step 0.

\emph{B6: the limit in $\rho$.}  Choose a sequence $\rho_j\uparrow R$ with
$0<\rho_j<R$, for instance $\rho_j=R(1-2^{-j-1})$.  For each $j$, (B6) holds for
every $x\in B_{\rho_j/2}$ outside a null set $N_j$; put $N=\bigcup_jN_j$, again a
null set.  Let $x\in B_{R/2}\setminus N$.  Since $|x-x_0|<R/2$ and
$\rho_j\uparrow R$, we have $|x-x_0|<\rho_j/2$ for all large $j$, so (B6) holds
at $x$ with $\rho=\rho_j$ for all large $j$:
\[
\bigl(0+|g(x)|^2\bigr)^{\frac p2}
\le\frac{C}{\rho_j^n}\int_{B_R}\bigl(0+|g|^2\bigr)^{\frac p2}dx .
\]
Letting $j\to\infty$ and using $\rho_j\to R$ together with the continuity of
$\rho\mapsto\rho^{-n}$ at $R>0$, the right-hand side converges to
$\frac C{R^n}\int_{B_R}(0+|g|^2)^{p/2}dx$, and therefore
\[
\bigl(\varepsilon+|g(x)|^2\bigr)^{\frac p2}
\le\frac{C}{R^n}\int_{B_R}\bigl(\varepsilon+|g|^2\bigr)^{\frac p2}dx
\qquad\text{for almost every }x\in B_{R/2},
\]
since $\varepsilon=0$.  Together with Step 2 this proves the theorem in both
cases, with the single constant $C$ of Step 0.

\emph{Remark: the thesis' homogeneity argument, and why it is not needed here.}
The thesis proves the case $R=1$ and removes the restriction by the rescaling of
pp.~17--18: for $t>0$ put $B_t(h)=A(h/t)$ and $u_t(x)=u(tx)$.  Those identities
are correct, and the three the thesis asserts ``it is easy to check'' are these.
Differentiating, $D_hB_t(h)=\frac1tD_hA(h/t)$, so for $\xi\in\R^n$
\[
\bigl(D_hB_t(h)\xi,\xi\bigr)=\frac1t\bigl(D_hA(h/t)\xi,\xi\bigr)
\ \ge\ \frac{\gamma_0}t\Bigl(\varepsilon+\frac{|h|^2}{t^2}\Bigr)^{\frac p2-1}|\xi|^2
=\frac{\gamma_0}{t^{p-1}}\bigl(t^2\varepsilon+|h|^2\bigr)^{\frac p2-1}|\xi|^2 ,
\]
because
$\bigl(\varepsilon+|h|^2/t^2\bigr)^{p/2-1}
=\bigl(t^{-2}(t^2\varepsilon+|h|^2)\bigr)^{p/2-1}
=t^{2-p}\bigl(t^2\varepsilon+|h|^2\bigr)^{p/2-1}$ and
$t^{-1}t^{2-p}=t^{1-p}$; and the same computation applied to (Q2) of the thesis gives
$|D_{h_j}B_t^i(h)|\le\gamma_1t^{1-p}\bigl(t^2\varepsilon+|h|^2\bigr)^{p/2-1}$.
So $B_t$ satisfies (Q1) of the thesis--(Q2) of the thesis with
$\varepsilon^t=t^2\varepsilon$, $\gamma_0^t=\gamma_0/t^{p-1}$ and
$\gamma_1^t=\gamma_1/t^{p-1}$, and the ratio
$\gamma_1^t/\gamma_0^t=\gamma_1/\gamma_0$ is unchanged --- which is also where
\S3's convention that $\varepsilon$ may be arbitrarily large is consumed, since
$t^2\varepsilon$ exceeds $1$ for large $t$.  The proof above does not use any of
this: \noderef{CaccioppoliPowerWeight} already holds on an arbitrary open set
with a constant depending on neither the set nor a radius, and
\noderef{MoserLocalBoundedness} is stated at an arbitrary radius in
scale-invariant form, so the reverse-H\"older family is available on $B_R$
directly and the iteration runs there.
\end{proof}
